\subsection{\texorpdfstring{Criterios de exactitud para los funtores
$H_\bullet(P_\bullet\otimes_A M)$}{Criterios de exactitud para los funtores H.(P. tensor M)}}
\label{subsection:III.7.4}
\label{III.7.4}

\begin{env}[7.4.1]
\label{III.7.4.1}
Sean $A$ un anillo (no necesariamente conmutativo) y $P_\bullet$ un complejo
de $A$-módulos planos a derecha. Como el funtor
$M\mapsto P_k\otimes_A M$ es entonces exacto en $\cat{Ab}_A$ para todo $k$,
el $\partial$-funtor
\[
  T_\bullet(M)=H_\bullet(P_\bullet\otimes_A M)
  \tag{7.4.1.1}\label{III.7.4.1.1.eq}
\]
\oldpage[III]{186}
es un funtor homológico de $\cat{Ab}_A$ en $\cat{Ab}$,
evidentemente $A$-lineal cuando $A$ es conmutativo
\sref{III.7.1.2}, y conmuta con los límites inductivos.

Si $A$ es conmutativo, entonces, para toda $A$-álgebra $B$, el funtor
homológico $T_\bullet^{(B)}$ \sref{III.7.3.8} viene dado, por definición, por
\[
  T_\bullet^{(B)}(N)
  =
  H_\bullet(P_\bullet\otimes_A N_{[\rho]})
  \tag{7.4.1.2}\label{III.7.4.1.2.eq}
\]
donde $\rho:A\to B$ es el homomorfismo que define la estructura de álgebra de
$B$; como también se puede escribir
\[
  P_\bullet\otimes_A N_{[\rho]}
  =P_\bullet\otimes_A(B\otimes_B N)_{[\rho]}
  =(P_\bullet\otimes_A B)\otimes_B N,
\]
se ve que
\[
  T_\bullet^{(B)}(N)=H_\bullet(P'_\bullet\otimes_B N)
  \tag{7.4.1.3}\label{III.7.4.1.3.eq}
\]
para todo $B$-módulo $N$, siendo $P'_\bullet$ el complejo
$P_\bullet\otimes_A B$ de $B$-módulos planos
\sref[0]{0.6.2.1}.
\end{env}

\begin{proposition}[7.4.2]
\label{III.7.4.2}
Bajo las condiciones generales de \sref{III.7.4.1}, y para un entero dado
$p\in\bb Z$, las propiedades siguientes son equivalentes:
\begin{enumerate}
  \item[(a)] $T_p$ es exacto por la izquierda (o, lo que es lo mismo,
    $T_{p+1}$ es exacto por la derecha);
  \item[(b)]
    \[
      Z'_p(P_\bullet)
      =
      \Coker(P_{p+1}\longrightarrow P_p)
    \]
    es un $A$-módulo plano a derecha;
  \item[(c)] existe un complejo $P'_\bullet$ de $A$-módulos planos a derecha
    tal que la diferencial
    \[
      d'_{p+1}:P'_{p+1}\longrightarrow P'_p
    \]
    sea nula, y un isomorfismo de funtores homológicos de
    $H_\bullet(P_\bullet\otimes_A M)$ sobre
    $H_\bullet(P'_\bullet\otimes_A M)$.
\end{enumerate}
\end{proposition}

\begin{proof}
Por definición, se tiene una sucesión exacta funtorial en $M$
\[
  0\longrightarrow T_p(M)
  \longrightarrow Z'_p(P_\bullet\otimes M)
  \longrightarrow P_{p-1}\otimes M
\]
donde
\[
  Z'_p(P_\bullet\otimes M)
  =
  \Coker(P_{p+1}\otimes M\longrightarrow P_p\otimes M)
  =
  Z'_p(P_\bullet)\otimes M
\]
en virtud de la exactitud por la derecha del producto tensorial. Para todo
homomorfismo $f:M\to N$, se tiene, por tanto, un diagrama conmutativo
\[
\xymatrix@C=30pt@R=38pt{
  0\ar[r]
  &T_p(M)\ar[r]\ar[d]_u
  &Z'_p(P_\bullet)\otimes M\ar[r]\ar[d]_v
  &P_{p-1}\otimes M\ar[d]_w
  \\
  0\ar[r]
  &T_p(N)\ar[r]
  &Z'_p(P_\bullet)\otimes N\ar[r]
  &P_{p-1}\otimes N
}
\tag{7.4.2.1}\label{III.7.4.2.1.eq}
\]
cuyas filas son exactas. Si $f$ es un monomorfismo, lo mismo ocurre con $w$,
puesto que $P_{p-1}$ es plano; si $T_p$ es exacto por la izquierda, $u$ es
también un monomorfismo; se concluye que $v$ es un monomorfismo, lo cual
implica que $Z'_p(P_\bullet)$ es plano. Recíprocamente, si es así, $v$ es un
monomorfismo para todo monomorfismo $f:M\to N$, luego el diagrama
\eqref{III.7.4.2.1.eq} muestra que $u$ es un monomorfismo, y, por consiguiente,
$T_p$ (que ya es semiexacto) es exacto por la izquierda. Así, (a) y (b) son
equivalentes. Es inmediato que (c) implica (a), pues si
$d'_{p+1}:P'_{p+1}\to P'_p$ es nula, y
$0\to M'\to M\to M''\to0$ es una sucesión exacta de $A$-módulos, el
homomorfismo de borde en la sucesión exacta
\[
  H_{p+1}(P'_\bullet\otimes M'')
  \xrightarrow{\partial}H_p(P'_\bullet\otimes M')
  \longrightarrow H_p(P'_\bullet\otimes M)
\]
\oldpage[III]{187}
es nulo por definición \emph{(M, IV, 1)}, luego $T_p$ es exacto por la
izquierda. Demostremos recíprocamente que (b) implica (c). Si
\[
  Z_{p+1}(P_\bullet)=\Ker(P_{p+1}\longrightarrow P_p),
\]
se tiene una sucesión exacta
\[
  0\longrightarrow Z_{p+1}(P_\bullet)\longrightarrow P_{p+1}
  \longrightarrow Z'_p(P_\bullet)\longrightarrow0
\]
en la cual $P_{p+1}$ y $Z'_p(P_\bullet)$ son planos, luego
$Z_{p+1}(P_\bullet)$ es plano \sref[0]{0.6.1.2}. Tomaremos
\[
  P'_i=P_i\quad\text{para }i\ne p\text{ y }i\ne p+1,
  \qquad
  P'_p=Z'_p(P_\bullet),
  \qquad
  P'_{p+1}=Z_{p+1}(P_\bullet);
\]
como diferencial $d'_i:P'_i\to P'_{i-1}$, se tomará la del complejo
$P_\bullet$ para $i\ne p$ e $i\ne p+1$, $0$ para $i=p+1$ y, para $i=p$,
el homomorfismo $Z'_p(P_\bullet)\to P_{p-1}$ deducido de $d_p$ por paso al
cociente. Como los $P_i$ son planos, se tiene
\[
  Z'_i(P_\bullet\otimes M)=Z'_i(P_\bullet)\otimes M,
  \qquad
  Z_i(P_\bullet\otimes M)=Z_i(P_\bullet)\otimes M
\]
y
\[
  B_i(P_\bullet\otimes M)=B_i(P_\bullet)\otimes M
\]
(poniendo $B_i(P_\bullet)=\Im(P_{i+1}\to P_i)$); se deducen inmediatamente,
para todo $M$, isomorfismos funtoriales
\[
  H_i(P_\bullet\otimes M)\xrightarrow{\sim}H_i(P'_\bullet\otimes M)
\]
para todo $i$, y la comprobación de que se trata de un isomorfismo de
$\partial$-funtores resulta sin dificultad de la definición de $\partial$
\emph{(M, IV, 1)}.
\end{proof}

Obsérvese que las condiciones de \sref{III.7.4.2} implican también que
$B_p(P_\bullet)$ es plano, pues se tiene una sucesión exacta
\[
  0\longrightarrow B_p(P_\bullet)\longrightarrow P_p
  \longrightarrow Z'_p(P_\bullet)\longrightarrow0,
\]
en la cual $P_p$ y $Z'_p(P_\bullet)$ son planos
\sref[0]{0.6.1.2}.

\begin{corollary}[7.4.3]
\label{III.7.4.3}
Supongamos que $A$ es un anillo noetheriano regular de dimensión $1$ (es
decir, un producto de anillos de Dedekind \sref[0]{0.17.1.3} y
\sref[0]{0.17.3.7}, por ejemplo un anillo de ideales principales). Entonces,
para que $T_p$ sea exacto por la izquierda, es necesario y suficiente que
$T_p(A)$ sea un $A$-módulo plano. Para que $T_p$ sea exacto, es necesario y
suficiente que $T_p(A)$ y $T_{p-1}(A)$ sean $A$-módulos planos.
\end{corollary}

\begin{proof}
Recordemos que, para un módulo $M$ sobre un anillo de Dedekind, es equivalente
decir que $M$ es plano o que es libre de torsión \sref[0]{0.6.3.3} y
\sref[0]{0.6.3.4}; bajo las hipótesis de \sref{III.7.4.3}, todo submódulo de un
$A$-módulo plano es, por tanto, plano.

La segunda aserción de \sref{III.7.4.3} resulta de la primera, puesto que decir
que $T_p$ es exacto significa que $T_p$ y $T_{p-1}$ son exactos por la
izquierda. Para demostrar la primera aserción, observemos que se tiene una
sucesión exacta
\[
  0\longrightarrow H_p(P_\bullet)\longrightarrow Z'_p(P_\bullet)
  \longrightarrow B_{p-1}(P_\bullet)\longrightarrow0
\]
en la cual $B_{p-1}(P_\bullet)$ es un $A$-módulo plano, por ser un submódulo
del $A$-módulo plano $P_{p-1}$. Por tanto, es equivalente decir que
$H_p(P_\bullet)$ es plano o que $Z'_p(P_\bullet)$ es plano
\sref[0]{0.6.1.2}.
\end{proof}

Las aplicaciones más importantes de \sref{III.7.4.2} son las siguientes:

\begin{proposition}[7.4.4]
\label{III.7.4.4}
Sean $A$ un anillo noetheriano y $P_\bullet$ un complejo de $A$-módulos
planos; se supone, o bien que los $P_i$ son de tipo finito, o bien que los
$H_i(P_\bullet)$ son $A$-módulos de tipo finito y que existe $i_0$ tal que
$H_i(P_\bullet)=0$ para $i<i_0$. Sea $T$ el funtor homológico definido por
\eqref{III.7.4.1.1.eq}. Entonces, el conjunto $U$ de los
$y\in\Spec(A)$ tales que $(T_p)_y$ \sref{III.7.1.4} sea exacto por la derecha
(resp. exacto por la izquierda, exacto) es abierto en $\Spec(A)$.
\end{proposition}

\begin{proof}
En la segunda hipótesis sobre $P_\bullet$, se puede reemplazar $P_\bullet$ por
un complejo $P'_\bullet$ de
\oldpage[III]{188}
$A$-módulos libres de tipo finito para el cual el funtor
$H_\bullet(P'_\bullet\otimes_A M)$ es isomorfo (como $\partial$-funtor) a
$T_\bullet(M)$ \sref[0]{0.11.9.3}. Por tanto, siempre se puede reducir a la
primera hipótesis, y, en este caso, los $Z'_i(P_\bullet)$ son de tipo finito;
además, basta limitarse a demostrar las aserciones relativas a la exactitud por
la izquierda (cf. \sref{III.7.4.2}, (a)). Dicho esto, sea $x\in U$; como el
funtor $M\mapsto M_x$ es exacto, se tiene
\[
  (Z'_p(P_\bullet))_x=Z'_p((P_\bullet)_x),
\]
y (teniendo en cuenta \eqref{III.7.4.1.3.eq}) la hipótesis implica, en virtud
de \sref{III.7.4.2}, (b), que $(Z'_p(P_\bullet))_x$ es un $A_x$-módulo plano,
luego libre, pues es de tipo finito y $A_x$ es un anillo local noetheriano
(Bourbaki, \emph{Algèbre commutative}, cap.~II, \S~3,
n\textsuperscript{o}~2, cor.~2 de la prop.~5). Se deduce que existe $f\in A$
tal que $(Z'_p(P_\bullet))_f$ es libre sobre $A_f$ (Bourbaki,
\emph{Algèbre commutative}, cap.~II, \S~5, n\textsuperscript{o}~1, cor.~de la
prop.~2), y, \emph{a fortiori}, $(Z'_p(P_\bullet))_y$ es libre sobre $A_y$ para
todo $y\in D(f)$, lo cual termina la demostración en virtud de
\sref{III.7.4.2}, (b).
\end{proof}

\begin{corollary}[7.4.5]
\label{III.7.4.5}
Bajo las hipótesis de \sref{III.7.4.4}, supongamos además que $A$ es un
dominio entero.
Entonces, el conjunto $U$ de los $x\in\Spec(A)$ tales que $(T_p)_x$ sea
exacto es abierto no vacío.
\end{corollary}

\begin{proof}
En efecto, basta demostrar que $(T_p)_x$ es exacto para el punto genérico $x$
de $\Spec(A)$, lo cual es inmediato, pues $A_x$ es un cuerpo, y, por tanto,
todo funtor aditivo sobre $\cat{Ab}_{A_x}$ es exacto.
\end{proof}

\begin{proposition}[7.4.6]
\label{III.7.4.6}
Bajo las hipótesis generales de \sref{III.7.4.4}, las condiciones (a), (b) y
(c) de \sref{III.7.4.2} son también equivalentes a:
\begin{enumerate}
  \item[(d)] existe un $A$-módulo $Q$ y un isomorfismo funtorial
    \[
      T_p(M)\xrightarrow{\sim}\Hom_A(Q,M).
      \tag{7.4.6.1}\label{III.7.4.6.1.eq}
    \]
    Además, el $A$-módulo $Q$ queda determinado, salvo isomorfismo único, por
    esta propiedad, y es de tipo finito.
\end{enumerate}
\end{proposition}

\begin{proof}
La unicidad de $Q$ es un caso particular de la unicidad de un objeto
representante de un funtor representable \sref[0]{0.8.1.5}. Es claro que el
segundo miembro de \eqref{III.7.4.6.1.eq} es exacto por la izquierda.
Recíprocamente, para demostrar la existencia de $Q$, cuando $T_p$ es exacto
por la izquierda, primero se puede, como en \sref{III.7.4.4}, reducir al caso
en que los $P_i$ son planos y de tipo finito, luego (puesto que $A$ es
noetheriano) proyectivos de tipo finito (Bourbaki,
\emph{Algèbre commutative}, cap.~II, \S~5, n\textsuperscript{o}~2, cor.~del
teor.~1). El dual $\check P_i$ de $P_i$ es entonces también un $A$-módulo
proyectivo de tipo finito, $P_i$ es canónicamente isomorfo al dual de
$\check P_i$, y el homomorfismo canónico
\[
  P_i\otimes_A M\longrightarrow\Hom_A(\check P_i,M)
\]
es biyectivo (Bourbaki, \emph{Algèbre}, cap.~II, 3.\textsuperscript{a} ed.,
\S~4, n\textsuperscript{o}~2, prop.~2). Por otra parte, se sabe por
\sref{III.7.4.2}, (c), que se puede suponer que
$d_{p+1}:P_{p+1}\to P_p$ es nula, de modo que se tiene una sucesión exacta
\[
  0\longrightarrow T_p(M)\xrightarrow{u}P_p\otimes M
  \xrightarrow{v}P_{p-1}\otimes M
\]
donde $v=d_p\otimes1$. Pongamos entonces $Q'=\Ker(d_p)$, de modo que se tiene
la sucesión exacta
\[
  0\longrightarrow Q'\xrightarrow{w}P_p\xrightarrow{d_p}P_{p-1},
\]
de donde, por transposición, la sucesión exacta
\[
  \check P_{p-1}\xrightarrow{{}^t d_p}\check P_p
  \xrightarrow{{}^t w}\check Q'\longrightarrow0.
\]
Vamos a ver que
\[
  Q=\check Q'=\Coker({}^t d_p)
\]
responde a la cuestión. En efecto, se tiene la sucesión exacta
\[
  0\longrightarrow\Hom_A(Q,M)
  \longrightarrow\Hom_A(\check P_p,M)
  \xrightarrow{v'}\Hom_A(\check P_{p-1},M)
\]
\oldpage[III]{189}
donde $v'=\Hom({}^t d_p,1)$; al identificar canónicamente
$P_i\otimes M$ con $\Hom(\check P_i,M)$, $v'$ se identifica, por tanto, con
$v=d_p\otimes1$, y se obtiene así el isomorfismo funtorial
\[
  T_p(M)\xrightarrow{\sim}\Hom_A(Q,M)
\]
buscado. Además, $Q$, por ser un cociente de $\check P_p$, es de tipo finito.
\end{proof}

\begin{proposition}[7.4.7]
\label{III.7.4.7}
Supongamos verificadas las condiciones generales de \sref{III.7.4.4}. Entonces,
para todo $A$-módulo $M$ de tipo finito:
\begin{enumerate}
  \item[(i)] los $T_i(M)$ son $A$-módulos de tipo finito;
  \item[(ii)] para todo ideal $\mathfrak m$ de $A$, el homomorfismo canónico
    \[
      \widehat{T_i(M)}
      \longrightarrow
      \varprojlim_n T_i\bigl(M\otimes_A(A/\mathfrak m^{n+1})\bigr)
      \tag{7.4.7.1}\label{III.7.4.7.1.eq}
    \]
    (donde el primer miembro es el completado separado de $T_i(M)$ para la
    topología $\mathfrak m$-ádica) es biyectivo.
\end{enumerate}
\end{proposition}

\begin{proof}
Como en \sref{III.7.4.4}, se reduce primero al caso en que los $P_i$ son de
tipo finito; como $A$ es noetheriano, los submódulos de los
$P_i\otimes_A M$ son de tipo finito, de donde se deduce trivialmente la
aserción (i). En cuanto a la aserción (ii), resulta más generalmente del lema
siguiente:
\end{proof}

\begin{lemma}[7.4.7.2]
\label{III.7.4.7.2}
Sean $A$ un anillo noetheriano y $u:E\to F$ un homomorfismo de $A$-módulos de
tipo finito. Para todo $A$-módulo $M$ de tipo finito, pongamos
\[
  K(M)=\Ker(u\otimes1_M),
  \qquad
  C(M)=\Coker(u\otimes1_M);
\]
entonces, los homomorfismos canónicos
\[
  \widehat{K(M)}\longrightarrow\varprojlim_n K(M_n),
  \qquad
  \widehat{C(M)}\longrightarrow\varprojlim_n C(M_n)
  \tag{7.4.7.3}\label{III.7.4.7.3.eq}
\]
(donde se ha puesto
$M_n=M\otimes_A(A/\mathfrak m^{n+1})=M/\mathfrak m^{n+1}M$) son biyectivos
para todo ideal $\mathfrak m$ de $A$.
\end{lemma}

\begin{proof}
Como $E\otimes M$ y $F\otimes M$ son de tipo finito, y el funtor
$M\mapsto\widehat M$ es exacto en la categoría de los $A$-módulos de tipo
finito \sref[0]{0.7.3.3}, $\widehat{K(M)}$ y $\widehat{C(M)}$ son,
respectivamente, el núcleo y el conúcleo de
\[
  \widehat{(u\otimes1)}:
  \widehat{(E\otimes M)}\longrightarrow\widehat{(F\otimes M)}.
\]
La exactitud por la izquierda del funtor $\varprojlim$ muestra, por tanto, que
$\widehat{K(M)}=\varprojlim_n K(M_n)$; por otra parte, la exactitud por la
derecha del producto tensorial prueba que
$C(M_n)=C(M)\otimes_A(A/\mathfrak m^{n+1})$, luego
$\widehat{C(M)}=\varprojlim_n C(M_n)$ por definición.
\end{proof}

\begin{remark}[7.4.8]
\label{III.7.4.8}
Teniendo en cuenta \sref{III.6.10.5} y \sref{III.6.10.6}, se ve que, bajo una
hipótesis suplementaria de planitud, \sref{III.7.4.7} vuelve a dar el hecho de
que (4.1.7.1) es un isomorfismo, es decir, lo esencial del «primer teorema de
comparación» para los morfismos propios; además, el enunciado se aplica ya no
solo a un $\sh{O}_X$-módulo coherente, sino a un complejo de tales módulos.
Sería interesante obtener un enunciado que comprendiese a la vez
\sref{III.7.4.7} y (4.1.7.1) como casos particulares. Obsérvese que, cuando los
$P_i$ son de tipo finito, la demostración de \sref{III.7.4.7} no utiliza el
hecho de que sean módulos planos; convendría examinar si la conclusión de
\sref{III.7.4.7} sigue siendo válida cuando no se supone que los $P_i$ sean
planos ni de tipo finito, sino que se supone que los $H_i(P_\bullet)$ son de
tipo finito para todo $i$ y nulos para $i<i_0$. ¿Es posible entonces reemplazar
$P_\bullet$ por un complejo $P'_\bullet$ de $A$-módulos de tipo finito tal que
los funtores $H_\bullet(P_\bullet\otimes M)$ y
$H_\bullet(P'_\bullet\otimes M)$ (que ya no son homológicos) sigan siendo
isomorfos?
\end{remark}
